\section{Testing expansion and finding sparse cuts in sparsifiers}
\label{sec:balsparse}
The key ingredient to our ED algorithms is that sparsifiers as per \cref{def:spars} serve as good proxies to the subgraphs of $G$ for any subroutine that either asserts expansion or outputs an approximate most-balanced sparse cut. We will in fact show that the result of running such subroutine on a boundary-linked subgraph of $G$ is almost the same as if we run it on the corresponding subgraph of a sparsifier.

Before we proceed further, we need to define what we demand exactly from the approximate most-balanced sparse cut subroutine. Loosely speaking, it is a procedure that takes as input a graph and a sparsity parameter, and either asserts that the graph is an expander, or outputs a sparse cut. In particular, the returned cut should be not too small compared to the largest cut among those certifying that the input is not an expander. For technical reasons, we ask that the algorithm also returns an approximation to the volume of the returned cut. Inspired by the definition of~\cite{streamexpdec}, we have the following definition, which captures what the subroutine should return.

\begin{definition}[Balanced sparse cut witness]
	\label{def:bscw}
	Let $\psi, \xi \in [0,1]$ and $\alpha,\lambda \ge 1$, and let $H=(U,E',w)$ be a weighted graph with self-loops. A $(\psi,\xi,\alpha,\lambda)$-balanced sparse cut witness of $H$ (for short, $(\psi,\alpha,\lambda,\xi)$-BSCW) is an element $\omega$ of $\{\bot\} \cup (\{R\,:\,\emptyset \neq R \subsetneq U\} \times \mathbb{R}_{>0})$ such that:
	\begin{enumerate}
		\item \label{item:expr} if $\omega = \bot$, then $H$ is a $\psi$-expander;
		\item if $\omega = (R,\nu)$, then:
		\begin{enumerate}
			\item \label{item:sparser} $\Phi_{H}(R) < \alpha \cdot \psi $,
			\item \label{item:balr} for every other cut $\emptyset \neq T \subsetneq U$ with $\Phi_H(T) <  \psi$ and $\vol_H(T) \le \vol_H(U \setminus T)$ one has $\vol_H(T) \le \lambda\cdot \vol_H(R)$,
			\item \label{item:smallsider} $\vol_H(R) \le (1+\xi) \vol_H(U \setminus R)$,
			\item \label{item:goodestimr} $ (1-\xi)  \cdot \unvolc{R}{U} \le \nu \le  (1+\xi)  \cdot \unvolc{R}{U}$.
		\end{enumerate}
	\end{enumerate}
\end{definition}
\noindent
A balanced sparse cut subroutine is then simply an algorithm returning a balanced sparse cut witness. For convenience of later discussions, we also have the following definition.

\begin{definition}[Balanced sparse cut algorithm]
	\label{def:bsca}
	Let $\alpha,\lambda \ge 1$. An algorithm $\bsc$ is an $(\alpha, \lambda)$-balanced sparse cut algorithm (for short, $(\alpha,\lambda)$-BSCA) if, taken as input a weighted graph with self-loops $H=(U,E',w)$ and a parameter $\psi \in (0,1)$, $\bsc(H,\psi)$ outputs a $(\psi,\alpha,\lambda,0)$-BSCW of $H$.
\end{definition}
\noindent
Ultimately, we will use either an exponential time brute force BSCA, or a fast BSCA with worse parameters, leading to our exponential and polynomial time ED algorithms, respectively.

Using these definitions, the main result of this section is the following: one can almost seamlessly run a BSCA on the sparsifier, without heavily deteriorating the quality of the BSCW hence obtained.
\begin{restatable}[Proxying lemma]{lemma}{reduction}
	\label{lem:reduction}
	Let $b, \phi, c,\delta \in (0,1)$ such that $c \leq 1/30$ is a constant, $\phi < b$, $b \le c$, $\delta \le c^2b/ \log n$, and  let $\bsc$ be an $(\alpha,\lambda)$-BSCA with $ \alpha \le \frac{1}{2b}$. Fix a cluster $U \subseteq V$ and also let $H=(V,E',w)$ such that $H[U] \approx_\delta G[U]$. Let $\omega$ be the output of $\bsc(H[U]^\circ, (1+\frac{1}{2\log n}) \cdot \phi)$. Then, $\omega$ is a  $(\phi,(1+\frac{1}{\log n})\alpha,(1+c)\lambda,c)$-BSCW of~$G[U]^\circ$.
\end{restatable}
\noindent
In \cref{subsec:boundlink} and \cref{subsec:preservingcuts}  we show properties that will let us compare cuts and volumes in $H[U]^\circ$ and $G[U]^\circ$. In \cref{subsec:proofproxying} we use these tools to prove \cref{lem:reduction}.

\subsection{Characterizing boundary-linked cuts}
\label{subsec:boundlink}
Structurally, we will divide the cuts of a subgraph $G[U]^\circ$ into two groups: those those that are boundary-linked and those that are not. This concept\footnote{The notion of boundary-linked cuts was already used in previous works, but not in the context of sparsification and streaming~\cite{boundlink,mincutjli}}., which we introduce next, is key to our goal: it captures the cuts that we are able to sparsify to within multiplicative error. Loosely speaking, a cut of $G[U]^\circ$ is boundary-linked if its local cut is a decent fraction of its global cut.

\begin{definition}[Boundary-linked cuts]
	\label{def:boundlinkcuts}
	Let $\beta \in (0,1)$ be a parameter, let $U \subseteq V$ be a cluster, and let $S \subseteq U$ be a cut. We say that $S$ is $\beta$-boundary-linked in $U$ if
	\begin{equation*}
		\uncut{S}{U} \ge \beta \cdot \min\{\uncutg{S}, \uncutg{(U \setminus S)}\} \, .
	\end{equation*}
\end{definition}
\noindent
We shall then group the cuts of $G[U]^\circ$ into boundary-linked ones and those that are not. As shown by the two lemmas below, each of these groups has a useful properties for our goal: for a correct set of parameters, boundary-linked cuts are very well approximated by a sparsifier, while those that are not boundary-linked must be sparse.

\begin{lemma}[Characterization of boundary-linked cuts]
	\label{lem:blinkcuts}
	Let $b, \phi, c,\delta \in (0,1)$ such that $c \leq 1/2$ is a constant, $\phi \le b$, $b \le c$, and $\delta \le c^2b/ \log n$.  Fix a cluster $U \subseteq V$ and also let $H=(V,E',w)$ such that $H[U] \approx_\delta G[U]$. Then, the following hold:
	\begin{enumerate}
		\item \label{charact:nonblink} for all non $b/2$-boundary-linked cuts $\emptyset \neq S \subsetneq U$ we have $\unspc{S}{U} < \phi$, and
		\item \label{charact:blink} for all $b/2$-boundary-linked cuts $S \subseteq U$ we have $(1-\frac{c}{\log n})\cdot \uncut{S}{U} \le \wcut{S}{U} \le 	(1+\frac{c}{\log n})\cdot \uncut{S}{U}$.
	\end{enumerate}
\end{lemma}
\begin{proof}
	To show~\eqref{charact:nonblink}, let $\emptyset \neq S \subsetneq U$ be a cut with $\unspc{S}{U} \ge \phi$. By definition this means that $\uncut{S}{U} \ge \phi \min \{\unvolc{S}{U},\unvolc{U \setminus S}{S}\}$. Since $\uncut{S}{U}=\uncut{(U\setminus S)}{U}$, we can in fact assume without loss of generality that $\uncut{S}{U} \ge \phi \unvolc{S}{U}$. Recalling how these volumes are defined, we have that every vertex in $u \in S$ has a contribution of at least $b\phi \cdot \unbor{\{u\}}{U}$ to $\unvolc{S}{U}$, so $\unvolc{S}{U} \ge b/\phi \cdot \unbor{S}{U}$. Hence, $\uncut{S}{U} \ge \phi \unvolc{S}{U}$ yields $\uncut{S}{U} \ge b \cdot \unbor{S}{U}$. Finally, we observe that the global cut of $S$, i.e.~$\uncutg{S}$, can be decomposed as the sum of $\uncut{S}{U}$ and $\unbor{S}{U}$. By virtue of this decomposition, from $\uncut{S}{U} \ge b \cdot \unbor{S}{U}$ we get $\uncut{S}{U} \ge b \cdot (\uncutg{S}-\uncut{S}{U})$. Rearranging the terms, and since we assume that~$b$ is bounded by $c \le 1/2$, we conclude $\uncut{S}{U} \ge b/2 \cdot \uncutg{S}$, i.e. $S$ is $b/2$-boundary-linked.
	
	To show~\eqref{charact:blink}, from definition of boundary-linked we know $\uncut{S}{U} \ge b/2 \cdot \min\{\uncutg{S}, \uncutg{(U \setminus S)}\} $, while from the assumption that $H[U] \approx_\delta G[U]$ we know $|\uncut{S}{U} - \wcut{S}{U} | \le \delta \cdot \uncutg{S}$. Again, without loss of generality we can assume $\uncut{S}{U} \ge b/2 \cdot \uncutg{S}$, since $\uncut{S}{U}=\uncut{(U\setminus S)}{U}$ and $\wcut{S}{U}=\wcut{(U\setminus S)}{U}$. Therefore, the sparsification error $|\uncut{S}{U} - \wcut{S}{U} | $ can actually be bounded by $\delta \cdot \uncutg{S} \le c^2b/ \log n \cdot 2/b \cdot \uncut{S}{U} \le  c/\log n \cdot \uncut{S}{U}$ (since $c\le 1/2$).
\end{proof}

\subsection{Preserving volumes and sparsities}
\label{subsec:preservingcuts}
 \cref{lem:blinkcuts} says that for any cut $S$ in $U$, either $S$ is sparse, or the sparsifier gives a multiplicative approximation to its local cut. Intuitively, for the latter cuts, one would expect the sparsifier to preserve their sparsity as well. However, this is not trivially true, because the sparsity depends on both the local cut and the volume in the boundary-linked subgraph: while we have a multiplicative preservation of boundary-linked cuts, we may not have the same guarantee for their volume. To circumvent this issue, we consider the case of sparse and expanding cuts separately: the former have in fact their volumes preserved within a relative error, while for the latter we get that their volume in the sparsifier is still not too large. We combine these observations by treating separately boundary-linked and non boundary-linked cuts, and conclude that the distinction between sparse and expanding is very accurately preserved in a sparsifier. We begin with analysing the volume of sparse cuts in the sparsifier.

\begin{lemma}[Multiplicative error for volumes of sparse cuts]
	\label{lem:volumesparsecuts}
	Let $b, \phi,\hat{\phi}, c,\delta \in (0,1)$ such that $c \leq 1/5$ is a constant, $\phi < b$, $\delta \le c^2 / \log n$, and \smash{$\hat{\phi} \le \phi/b$}. Fix a cluster $U \subseteq V$ and also let $H=(V,E',w)$ such that $H[U] \approx_\delta G[U]$. Then, for all cuts $\emptyset \neq S \subsetneq U$ with \smash{$\unspc{S}{U} < \hat{\phi}$} we have
	\begin{equation*}
		\left(1-\frac{c}{\log n}\right) \cdot \unvolc{S}{U} \le \wvolc{S}{U} \le \left(1+\frac{c}{\log n}\right) \cdot \unvolc{S}{U} \, .
	\end{equation*}
\end{lemma}
\begin{proof}
	By definition of $\wvolc{S}{U}$, the error in estimating the volume can be written as
	\begin{equation}
		\left|\wvolc{S}{U}- \unvolc{S}{U}\right| \le \left|\wvol{S}- \vol_U(S)\right| + \frac{b-\phi}{\phi}\left|\wbor{S}{U} - \unbor{S}{U}\right| \, . \label{eq:triangineqselfloops}
	\end{equation}
	The first term is easily upper-bounded by $\delta \cdot \vol(S)$, because of the global cut preservation property of $H[U] \approx_\delta G[U]$ (see \cref{def:spars}) applied to singleton cuts. For the second term, we can further split it as
	\begin{equation}
		\left|\wbor{S}{U} - \unbor{S}{U}\right| \le \left|\wcutg{S}-\uncutg{S}\right|+\left|\wcut{S}{U}-\uncut{S}{U}\right| \, , \label{eq:trianginequality}
	\end{equation}
We now upper-bound the right hand side using again the cut preservation properties from \cref{def:spars}: we employ the multiplicative error for the global cut $\uncutg{S}$, and the additive error for $\uncut{S}{U}$. Combining these with the assumption that $S$ is sparse, one gets
	\begin{align}
		\Big|\wcutg{S} - \uncutg{S}\Big| + \Big|\wcut{S}{U}-\uncut{S}{U}\Big| 
		& \le  2\delta\cdot \uncutg{S} \\
		& =   2\delta \cdot\big(\uncut{S}{U} + \unbor{S}{U}\big) \\
		  \text{ since $\unspc{S}{U} < \hat{\phi}$ } \quad & <  2\delta \cdot \hat{\phi} \cdot \vol_U^\circ(S) + 2\delta \cdot \unbor{S}{U}\\
	 \text{ since $\hat{\phi} \le \phi/b$ and $\vol_U^{\circ}(S) \ge \frac{\phi}{b}\unbor{S}{U}$ } \quad & \le 2\delta \cdot \frac{\phi}{b} \cdot \vol_U^\circ(S) + 2\delta \cdot \frac{\phi}{b} \cdot \vol_U^{\circ}(S)\\
		& \le 4\delta\cdot\frac{\phi}{b-\phi} \cdot \vol_U^{\circ}(S) \, , \label{eq:sumgloballocalerror}
	\end{align}
We now backtrack and plug~\eqref{eq:sumgloballocalerror} into~\eqref{eq:trianginequality}, and then~\eqref{eq:trianginequality} into~\eqref{eq:triangineqselfloops} to conclude the claim:
	\begin{align*}
		\left|\wvolc{S}{U}- \unvolc{S}{U}\right| & \le \left|\wvol{S}- \vol_U(S)\right| + \frac{b-\phi}{\phi}\left|\wbor{S}{U} - \unbor{S}{U}\right| \\
		& \le \delta \cdot \vol(S) + 4\delta \cdot \vol_U^{\circ}(S) \\
		& \le 5\delta \cdot \vol_U^{\circ}(S) \\
		& \le \frac{5c^2}{\log n}\cdot \unvolc{S}{U} \\
		& \le \frac{c}{\log n}\cdot \unvolc{S}{U} \, ,
	\end{align*}
since $c \le 1/5$.
\end{proof}
\noindent
Using this fact together with the cut preservation properties of a cluster sparsifier, we can already conclude that there cannot be cuts that look ``very'' expanding in $H[U]^\circ$ but are actually sparse in~$G[U]^\circ$. In particular, if the sparsity of a cut in the sparsifier is above $\phi$ by a factor that dominates the error of the sparsifier, then its sparsity must be at least $\phi$ in the original graph. We show this in \cref{lem:iffsparsity}.
\begin{lemma}[Very expanding cuts in a sparsifier are indeed expanding]
	\label{lem:iffsparsity}
	Let $b, \phi,\hat{\phi}, c,\delta \in (0,1)$ such that $c \leq 1/30$ is a constant, $\phi < b$, $b \le c$, $\delta \le c^2 b/ \log n$, and \smash{$\hat{\phi} \ge (1+\frac{1}{3\log n})\cdot \phi$}.  Fix a cluster $U \subseteq V$ and also let $H=(V,E',w)$ such that $H[U] \approx_\delta G[U]$. Then, for all cuts $\emptyset \neq S \subsetneq U$ with \smash{$\wspc{S}{U} \ge \hat{\phi}$} we have $\unspc{S}{U} \ge \phi$.
\end{lemma}
\begin{proof}
	For the sake of a contradiction, let $\emptyset \neq S \subsetneq U$ be a cut with \smash{$\wspc{S}{U} \ge \hat{\phi}$} such that $\unspc{S}{U} < \phi$.
	
	Consider first the case that $S$ is $b/2$-boundary-linked. Then we know from $H[U] \approx_\delta G[U]$ (see \cref{def:spars}) that $\wcut{S}{U} = (1 \pm c/\log n) \uncut{S}{U}$. Also, as we are assuming $\unspc{S}{U} < \phi$, from \cref{lem:volumesparsecuts} (our parameters verify stronger conditions than what is demanded by the lemma, so it applies) we know $\wvolc{S}{U}= (1 \pm c/\log n)\unvolc{S}{U}$. These two together give $\wspc{S}{U} < (1+10c/\log n)\phi < \hat{\phi}$ for small enough $c$, a contradiction.
	
	Next consider the case that $S$ is not $b/2$-boundary-linked. Then we do not have a multiplicative approximation to $\wcut{S}{U}$. However, $S$ is $\phi$-sparse in $G[U]^\circ$ so we still have $\wvolc{S}{U}= (1 \pm c/\log n)\unvolc{S}{U}$ and $\wvolc{U \setminus S}{U}= (1 \pm c/\log n)\unvolc{U \setminus S}{U}$ by Lemma~\ref{lem:volumesparsecuts}. Therefore, because $\hat{\phi} \ge (1+\frac{1}{3\log n}) \phi$ we get
	\begin{align*}
		\wcut{S}{U} & \ge \left(1-\frac{c}{\log n}\right) \left(1+\frac{1}{3\log n}\right)\cdot \phi \cdot \min\{\unvolc{S}{U}, \unvolc{U \setminus S}{U}\} \\
		& \ge \left(1+\frac{1}{10\log n}\right)\cdot \phi \cdot \min\{\unvolc{S}{U}, \unvolc{U \setminus S}{U}\} \, ,
	\end{align*}
	and at the same time the additive error approximation from \cref{def:spars} is enough to yield
	\begin{equation*}
		\wcut{S}{U} \le \uncut{S}{U} + \delta \cdot \uncutg{S} < \phi \cdot \min\{\unvolc{S}{U}, \unvolc{U \setminus S}{U}\}  + \delta\cdot \uncutg{S} \, .
	\end{equation*}
	Therefore,
	\begin{align*}
		\uncutg{S} & \ge \frac{1}{\delta \cdot10 \log n} \cdot \phi \cdot \min\{\unvolc{S}{U}, \unvolc{U \setminus S}{U}\} \\
		\text{ since $\delta \le c^2 b/ \log n$ } \quad & \ge \frac{1}{10c^2} \cdot \min \{\unbor{S}{U}, \unbor{U \setminus S}{U}\}\, .
	\end{align*}
	Now let us assume $\unbor{S}{U} \le \unbor{U \setminus S}{U}$ without loss of generality. Then the above is essentially saying that $\uncutg{S}$ is much larger than $\unbor{S}{U}$.	Recalling now that we can decompose $\uncutg{S}=\uncut{S}{U}+\unbor{S}{U}$, the above lower bound is actually saying that much of the advantage that~$\uncutg{S}$ has over~$\unbor{S}{U}$ must come from $\uncut{S}{U}$. More precisely, since $c \le 1/30$ we get
	\begin{equation*}
		\uncut{S}{U} \ge\frac{1}{c} \cdot \uncutg{S} > \frac{b}{2} \cdot \uncutg{S} \, ,
	\end{equation*}
	which contradicts the assumption that $S$ is not $b/2$-boundary-linked.
\end{proof}
\noindent
Next, we handle the case of expanding cuts, and show that their volume does not get overshot too much. Even though this is a much weaker guarantee than the one we obtained in Lemma~\ref{lem:volumesparsecuts} for sparse cuts, it will  still be enough for our purposes: this will later be useful in concluding that their sparsity has not dropped by a lot.

\begin{lemma}[Upper bound for volumes of expanding cuts]
	\label{lem:testexpandingcuts}
	Let $b, \phi,\hat{\phi}, c,\delta \in (0,1)$ such that $c \leq 1/2$ is a constant, $\phi \le b$, $\delta \le c^2 / \log n$, and \smash{$\hat{\phi} \le \phi/b$}. Fix a cluster $U \subseteq V$ and also let $H=(V,E',w)$ such that $H[U] \approx_\delta G[U]$. Then, for all cuts $\emptyset \neq S \subsetneq U$ with \smash{$\Phi_U^\circ(S) \ge \hat{\phi}$} we have
	\begin{equation*}
		\min \{\wvolc{S}{U},\wvolc{U \setminus S}{U}\} \le \frac{1+\frac{c}{\log n}}{\hat{\phi}} \cdot \uncut{S}{U}\, .
	\end{equation*}
\end{lemma}
\begin{proof}
	Let $S^* \in \{S, U \setminus S\}$ be the side of the cut achieving the minimum $\min\{\unvolc{S}{U}, \unvolc{U \setminus S}{U}\}$. Even though $\wbor{S^*}{U}$ may not be a very good approximation of $\unbor{S^*}{U}$, we can upper-bound it by decomposing  $\wcutg{S^*}=\wcut{S}{U}+\wbor{S}{U}$ and using the local cut preservation property of $H[U] \approx_\delta G[U]$ (see \cref{def:spars}): we get
	\begin{align*}
		\wbor{S^*}{U} & = \wcutg{S^*} - \wcut{S^*}{U} \le  \unbor{S^*}{U}+\delta \cdot \uncutg{S^*} \, .
	\end{align*}
Using the above bound, together with the global approximation from \cref{def:spars} applied to singleton cuts, we can upper-bound $\wvolc{S^*}{U}$ as
	\begin{align*}
		\wvolc{S^*}{U} & = \wvol{S}+\frac{b-\phi}{\phi}\wbor{S}{U} \\
		& \le  (1+\delta) \cdot \vol(S^*)+\frac{b-\phi}{\phi} \unbor{S^*}{U} + \frac{b-\phi}{\phi}  \cdot \delta \cdot \uncutg{S^*}\\
		\text{ since $\uncutg{S^*}=\uncut{S}{U}+\unbor{S}{U}$} \quad & \le (1+\delta) \cdot \vol(S^*)+ (1+\delta) \cdot\frac{b-\phi}{\phi}  \cdot \unbor{S^*}{U}+ \frac{b-\phi}{\phi}  \cdot \delta \cdot \uncut{S^*}{U}\\
		\text{ by definition of $\unvolc{S}{U}$} \quad & \le (1+\delta) \cdot \vol_U^{\circ}(S^*)+ \frac{b-\phi}{\phi}  \cdot \delta \cdot \uncut{S^*}{U}\\
		\text{ since $\hat{\phi} \le \frac{\phi}{b}$} \quad & \le (1+\delta) \cdot  \unvolc{S^*}{U} + \delta \cdot  \frac{1}{\hat{\phi}} \cdot \uncut{S^*}{U} \, .
	\end{align*}
By our assumption that $\Phi_U^{b/\phi}(S) \ge \hat{\phi}$, we also know $\vol_U^{b/\phi}(S^*) \le \frac{1}{\hat{\phi}} \uncut{S}{U}$. We then conclude
\begin{align*}
	\wvolc{S^*}{U}
	& \le (1+2\delta) \cdot \frac{1}{\hat{\phi}} \uncut{S}{U} \le \left(1+\frac{c}{\log n}\right) \frac{1}{\hat{\phi}} \uncut{S}{U}\, ,
\end{align*}
since $c \le 1/2$.
\end{proof}
\noindent
Having a bound on the volumes of expanding cuts, allows to bound the sparsity of those cuts in the sparsifier (since all such cuts are boundary-linked, and hence the local cut is well preserved). This is almost the reverse of \cref{lem:iffsparsity}.

\begin{lemma}[Expanding cuts are almost expanding in the sparsifier]
	\label{lem:boundlinkexpsparse}
Let $b, \phi,\hat{\phi}, c,\delta \in (0,1)$ such that $c \leq 1/6$ is a constant, $\phi \le b$, $b \le c$, $\delta \le c^2 b/ \log n$, and \smash{$\phi \le \hat{\phi} \le \phi/b$}.  Fix a cluster $U \subseteq V$ and also let $H=(V,E',w)$ such that $H[U] \approx_\delta G[U]$. Then, for all cuts $\emptyset \neq S \subsetneq U$ with \smash{$\unspc{S}{U} \ge \hat{\phi}$} we have $\wspc{S}{U} \geq (1 - \frac{1}{3\log n}) \hat{\phi}$.
\end{lemma}
\begin{proof}
	First note that our parameter regime is strictly stronger than that of \cref{lem:blinkcuts} and \cref{lem:testexpandingcuts}. Hence, by \cref{lem:blinkcuts}, we know that $S$ is $b/2$-boundary-linked, since it is $\hat{\phi}$-expanding and $\hat{\phi} \ge \phi$. \cref{lem:blinkcuts} then also guarantees that $\wcut{S}{U} \ge (1-\frac{c}{\log n}) \uncut{S}{U}$. On the other hand, \cref{lem:testexpandingcuts} gives \smash{$	\min \{\wvolc{S}{U},\wvolc{U \setminus S}{U}\} \le (1+\frac{c}{\log n})/{\hat{\phi}} \cdot \uncut{S}{U}$}. Hence,
	\begin{equation*}
		\min \{\wvolc{S}{U},\wvolc{U \setminus S}{U}\} \le \frac{1+\frac{c}{\log n}}{1-\frac{c}{\log n}}\frac{1}{\hat{\phi}} \cdot \wcut{S}{U} \leq  \frac{1}{1-\frac{1}{3\log n}}\cdot \frac{1}{\hat{\phi}} \cdot \wcut{S}{U}  \, ,
	\end{equation*}
since $c \le 1/6$ is small enough.
\end{proof}

\subsection{Proving the proxying lemma}
\label{subsec:proofproxying}
From the previous section, we have that the distinction between sparse and expanding cuts of $G[U]^\circ$ is very accurately preserved by $H[U]^\circ$, and also the volume of sparse cuts of $G[U]^\circ$ is approximated to within a small relative error in $H[U]^\circ$. As these are exactly the quantities that a BSCA is interested in, we can then prove the main lemma of the section, restated here for convenience of the reader.

\reduction*

\begin{proof}[Proof]
	By definition of $\bsc$, we know that $\omega$ is a $((1+\frac{1}{2\log n}) \cdot \phi, \alpha,\lambda,0)$-BSCW of~$H[U]^\circ$, i.e.
	\begin{enumerate}
		\item \label{itdef:exp2} if $\omega=\bot$, then $H[U]^\circ$ is a $(1+\frac{1}{2\log n}) \cdot \phi$-expander;
		\item if $\omega=(R,\nu)$, then:
		\begin{enumerate}
			\item \label{itdef:sparse2} $\wspc{R}{U} < (1+\frac{1}{2\log n}) \cdot \alpha \cdot \phi $,
			\item \label{itdef:smallside2} $\wvolc{R}{U} \le \wvolc{U \setminus R}{U} $,
			\item \label{itdef:bal2} for every other cut $\emptyset \neq T \subsetneq U$ with $\wspc{R}{U} <(1+\frac{1}{2\log n}) \cdot \phi$ and $\wvolc{T}{U} \le \wvolc{U \setminus T}{U}$ one has $\wvolc{T}{U} \le B \wvolc{R}{U} $:
			\item \label{itdef:goodestim2} $\nu =\wvolc{R}{U}$.
		\end{enumerate}
	\end{enumerate}
Our goal is to translate these properties to $G[U]^\circ$ by showing that $\omega$ is a  $(\phi,(1+\frac{1}{\log n})\alpha,(1+c)\lambda,c)$-BSCW of~$G[U]^\circ$, i.e.
\begin{enumerate}
	\item \label{itdef:exp3} if $\omega=\bot$, then $G[U]^\circ$ is a $ \phi$-expander;
	\item if $\omega=(R,\nu)$, then:
	\begin{enumerate}
		\item \label{itdef:sparse3} $\unspc{R}{U} < (1+\frac{1}{\log n}) \cdot\alpha \cdot  \phi $,
		\item \label{itdef:smallside3} $\unvolc{R}{U} \le (1+c)\unvolc{U \setminus R}{U} $,
		\item \label{itdef:bal3} for every other cut $\emptyset \neq T \subsetneq U$ with $\unspc{R}{U} <\phi$ and $\unvolc{T}{U} \le \unvolc{U \setminus T}{U}$ one has $\unvolc{T}{U} \le (1+c)B \unvolc{R}{U} $:
		\item \label{itdef:goodestim3} $(1-c)\cdot\unvolc{R}{U}\le \nu \le (1+c)\cdot\unvolc{R}{U}$.
	\end{enumerate}
\end{enumerate}

	\paragraph{Correctness of the case $\omega = \bot$.} The idea is that if the cluster $U$ is non-trivially more than a $\phi$-expander in the sparsifier, then it should be at least a $\phi$-expander in $G$. Let us say that $\omega=\bot$, i.e. $\wspc{S}{U} \ge  (1+\frac{1}{2\log n})\cdot \phi $ for all $\emptyset \neq S \subsetneq U$. Let \smash{$\hat{\phi}= (1+\frac{1}{2\log n}) \cdot \phi$}, so that we are in the parameter regime of \cref{lem:iffsparsity}. Then, this lemma ensures that $\unspc{S}{U} \ge \phi$ for all $\emptyset \neq S \subsetneq U$. In other words, $G[U]^\circ$ is a $\phi$-expander.
	
	\paragraph{Correctness of the case $\omega=(R,\nu)$.}
Assume $\omega=(R,\nu)$. We need to show the four properties of a BSCW for~$G[U]^\circ$.
	\begin{itemize}
		\item \textit{Property~\eqref{itdef:sparse3}}. Intuitively, $R$ cannot be much more expanding in $G$ than what it looks like in $H$, so if it is sparse in $H$ it is basically just as sparse in $G$. To prove this formally, it is convenient to distinguish the case of this cut being boundary-linked or not.
		
		If $R$ is not $b/2$-boundary-linked, then \cref{lem:blinkcuts} implies $\unspc{R}{U} < \phi$. The property is  then already proved in this case. 
		
		If it is $b/2$-boundary-linked, suppose for the sake of a contradiction that $R$ has {$\unspc{R}{U} \ge (1+{1}/{\log n})\alpha \cdot \phi$}.  Let \smash{$\hat{\phi}=(1+\frac{1}{\log n})\alpha \cdot \phi$}. This value of \smash{$\hat{\phi}$} fulfils the condition that \smash{$\hat{\phi} \le 1/b$} from \cref{lem:boundlinkexpsparse}, since $\alpha\le 1/(2b)$ by assumption. We then meet the requirements to apply \cref{lem:boundlinkexpsparse} (all other parameters also meet its preconditions, as our parameter regime for $b,\phi,c,\phi$ is no weaker than its), so
		
		\[
			\wspc{R}{U} \geq \left(1-\frac{1}{3\log n}\right)\cdot \left(1+\frac{1}{\log n}\right)\cdot\alpha \cdot \phi \geq \left(1+\frac{1}{2\log n}\right) \alpha \cdot \phi,
		\]
		which contradicts the definition of $\bsc$. We have then showed property~\eqref{itdef:sparse3}.
		
		
		
		\item \textit{Properties~\eqref{itdef:smallside3} and~\eqref{itdef:goodestim3}}. From the previous point we know that $\unspc{R}{U} <(1+\frac{1}{\log n})\alpha \cdot \phi$. Since our parameter regime is only stronger than that of \cref{lem:volumesparsecuts} if we set \smash{$\hat{\phi}=(1+\frac{1}{\log n})\alpha \cdot \phi$}, we get $\wvolc{R}{U}= (1 \pm c/\log n)\unvolc{R}{U}$ and $\wvolc{U \setminus R}{U}= (1 \pm c/\log n)\unvolc{U \setminus R}{U}$. Recalling that $\wvolc{R}{U} \le \wvolc{U\setminus R}{U}$ (see~\eqref{itdef:smallside2} for $H[U]^\circ$), this yields
		\begin{equation*}
			\unvolc{R}{U} \le \frac{1}{1-\frac{c}{\log n}} \wvolc{R}{U} \le  \frac{1}{1-\frac{c}{\log n}}\wvolc{U\setminus R}{U} \le \frac{1+\frac{c}{\log n}}{1-\frac{c}{\log n}}\unvolc{U\setminus R}{U} \, ,
		\end{equation*}
		so $\unvolc{R}{U} \le (1+c)\unvolc{U\setminus R}{U}$. Recalling that $\nu = \wvolc{R}{U}$ (see~\eqref{itdef:goodestim2} for $H[U]^\circ$), one also has
		\begin{equation*}
			\left|\nu-\unvolc{R}{U}\right|=\left|\wvolc{R}{U}-\unvolc{R}{U}\right|\le c/\log n \cdot \unvolc{R}{U} \le c \unvolc{R}{U} \, .
		\end{equation*}
		
		
		
		\item \textit{Property~\eqref{itdef:bal3}}. This property aims to bound the volume of $R$ against the volume of every other $\phi$-sparse cut. We want to exploit the bound that $\bsc$ gives for these cuts in $H$, i.e. that for every cut $\emptyset \neq T \subsetneq U$ with $\wspc{R}{U} <  (1+\frac{1}{2\log n})  \phi$ and $\wvolc{T}{U} \le \wvolc{U \setminus T}{U}$ one has $\wvolc{T}{U} \le \lambda \wvolc{R}{U} $ (see~\eqref{itdef:bal2} for $H[U]^\circ$). Now one can see that if we prove all such $T$ to be also $\phi$-sparse in $G[U]^\circ$, then we could conclude simply using the fact that the volumes of $H[U]^\circ$ are good proxies for those in $G[U]^\circ$. We now implement this strategy.
		
		By virtue of property~\eqref{itdef:sparse3}, we have that $\unspc{R}{U}<(1+\frac{1}{\log n})\alpha \cdot \phi$, and from an application of \cref{lem:volumesparsecuts} we have $\wvolc{R}{U}= (1 \pm c/\log n)\unvolc{R}{U}$ (again, the preconditions of the lemma are met if we take \smash{$\hat{\phi}=(1+\frac{1}{\log n})\alpha \cdot \phi$}).
		
		Consider now any other cut $\emptyset \neq T \subsetneq U$ with $\unspc{T}{U} < \phi$ and $\unvolc{T}{U} \le \unvolc{U \setminus T}{U}$. Note that for any such $T$ we have $\wvolc{T}{U}= (1 \pm c/\log n)\unvolc{T}{U}$ and $\wvolc{U \setminus T}{U}= (1 \pm c/\log n)\unvolc{U \setminus T}{U}$, again by applying \cref{lem:volumesparsecuts} (setting $\hat{\phi} =\phi \le \phi/b$, and all other parameters satisfy the needed preconditions).
		
		Moreover, any such $T$ is also $((1+\frac{1}{2\log n})\ \cdot \phi)$-sparse in $H[U]^\circ$, by applying Lemma~\ref{lem:iffsparsity} with parameter \smash{$\hat{\phi}=(1+\frac{1}{2\log n}) \phi \ge (1+\frac{1}{3\log n}) \phi$}. From~\eqref{itdef:bal2} for $H[U]^\circ$, we then know
		\begin{equation*}
			\min\{\wvolc{T}{U},\wvolc{U \setminus T}{U}\} \le \lambda \wvolc{R}{U} \, .
		\end{equation*}
		Recall that $\wvolc{R}{U}$, $\wvolc{T}{U}$, $\wvolc{U \setminus T}{U}$ all are within a $(1 \pm c/\log n)$ multiplicative error of $\unvolc{R}{U}$, $\unvolc{T}{U}$, $\unvolc{U \setminus T}{U}$ respectively. Hence, $\min\{\wvolc{T}{U},\wvolc{U \setminus T}{U}\} \le \lambda \wvolc{R}{U}$ translates to
		\begin{equation*}
			\unvolc{T}{U} = \min\{\unvolc{T}{U},\unvolc{U \setminus T}{U}\} \le (1+10c/\log n) \lambda\cdot \unvolc{R}{U} \le (1+c) \lambda\cdot \unvolc{R}{U} \, .
		\end{equation*}
	\end{itemize}
	\noindent
	We have then proved~\eqref{itdef:bal3}.
	\end{proof}